\chapter{Laumon 层}\label{ch:4}\index{Laumon 层|hyperpage}
% 原PDF第26页中部
设 $X$ 是 $k$ 上光滑射影连通曲线，$E$ 是 $X$ 上秩 $n$ 的局部系统。局部描述是在 étale 局部意义下的，因此下文 $\A^1$ 上的例子只是局部模型。目标是在 $X$ 上挠层的叠 $\Tor$ 上构造反常层 $\cL_E$。这是 Laumon 构造 $\GL_n$ 自守层时用来替代 Whittaker 函数的层论输入\cite{FGKV98}；这一构造的抛物变体与分歧变体见文献\cite{Hei04}。
\section{来自 Whittaker 与 Satake 的动机}\label{sec:4.1}
Whittaker 讲授（见\chref{3}）中理想的图景为
\origdiagram{p026-diagram.pdf}{0.5829}
这里 $\delta_E$ 是局部系统 $E$ 处的摩天层。应当把希望得到的 Whittaker 对象理解为
\[W_E\simeq\delta_E*W_0,\]
其中 $W_0$ 是真空 Whittaker 层。将 $\checkG$ 的表示 $V$ 作用于整个 $\checkG$-局部系统 $E$，得到 $X$ 上的局部系统 $V_E$，其在 $x$ 处的纤维为 $V_{E,x}$。对 $G=\GL_n$，这是对秩 $n$ 局部系统 $E$ 的通常操作。Laumon 层是在 $G=\GL_n$ 的正 Hecke 轨迹上对这一原则的具体实现。

\begin{transnote}{含糊点拨}
卷积 $*$ 是 Whittaker 范畴 $\operatorname{Wh}(\dot{\Gr}_{G,x})$ 上沿 Hecke 对应所作的 $!$-卷积：两个对象卷起来，就是先用一次修改把其中一个搬运过去、再叠上另一个，等价于把「作用」这一操作并入范畴。$W_0$ 是真空 Whittaker 层，即对应平凡 $\checkG$-表示（真空向量）的那个对象；$\delta_E$ 是 $E$ 处的摩天层，作用是把 $\checkG$ 的表示 $V$ 的实现「钉」在局部系统 $E$ 上。于是 $\delta_E*W_0$ 恰好是「把 $V_E$ 作用于真空」所得的 Whittaker 对象，也就是几何 Satake 的 $\operatorname{Sat}_{V_E}$ 在 Whittaker 一侧的对应物——这就是把 $W_E$ 写成这种卷积形式的意思。
[依据：原书 p.26–27 \S\ref{sec:4.1} 的图景；并见 [Gai18] 对 Whittaker 范畴上卷积的构造与 [FGKV98]（JAMS 11 (1998), 451–484）的 Whittaker 与 Satake 对应。]
\end{transnote}$\dot{\Gr}_{G,x}$ 中的上点是用于 Whittaker 层的同一个 $\rho$ 扭曲；下文的球面 Satake 输入写在普通仿射 Grassmann 簇上。

局部输入是几何 Satake：
\[\operatorname{Perv}_{G(\cO_x)}(\Gr_{G,x})\simeq\Rep(\checkG),\qquad\operatorname{IC}_{\overline{\Gr_{G,x}^\lambda}}\longleftrightarrow V^\lambda,\quad\lambda\in\Lambda^+.\]
采用 $\Gr_{G,x}\simeq G(K_x)/G(\cO_x)$ 及 $\Gr_{G,x}^\lambda=G(\cO_x)t_x^\lambda G(\cO_x)/G(\cO_x)$。

还存在几何 Casselman–Shalika 形式
\[\operatorname{Wh}(\dot{\Gr}_{G,x})\simeq\Rep(\checkG),\qquad W_\lambda\longleftrightarrow V^\lambda,\]
它与 Satake 作用相容。采用\chref{3}的 $\rho$ 扭曲记号，Whittaker 轨道为
\[S^\lambda:=\dot N(K_x)t_x^\lambda\dot G(\cO_x)/\dot G(\cO_x),\qquad\lambda\in\Lambda.\]
\srcpage{27}
$S^\lambda$ 上存在 $(\dot N(K_x),\Psi_x)$-等变层，当且仅当 $\lambda$ 优势；事实上，
\[\operatorname{Stab}_{\dot N(K_x)}(t_x^\lambda)\subset\ker\Psi_x\quad\Longleftrightarrow\quad\lambda\in\Lambda^+.\]
对这样的 $\lambda$，$S^\lambda$ 上相应的扭曲局部系统可洁净延拓到 $\dot{\Gr}_{G,x}$。一维模型是 Artin–Schreier 层：对 $j:\A^1\hookrightarrow\mathbb P^1$ 及 $i_\infty:\{\infty\}\hookrightarrow\mathbb P^1$，
\[i_\infty^*j_*\cL_\psi=0,\qquad\text{从而 }j_!\cL_\psi\simeq j_*\cL_\psi.\]
关键是 $j_*\cL_\psi$ 在边界点 $\infty$ 处的茎消失；对于零延拓，$i_\infty^!j_!$ 的消失则是恒然成立的。这是非退化 Whittaker 局部系统洁净延拓的一维模型。

\begin{transnote}{含糊点拨}
Artin–Schreier 层 $\cL_\psi$ 是 $\A^1$ 上由覆盖 $x\mapsto x^p-x$（$k$ 特征 $p$）所给的 $\mathbb Z/p$-局部系统在 $\psi$ 下的扭：它的单值性正是 $\psi$，因此在 $\infty$ 处有正则奇点。关键的一步是 $j_!$ 与 $j_*$ 只在边界处不同：对开浸入 $j$，锥 $\operatorname{Cone}(j_!\to j_*)$ 支撑在边界点 $\infty$ 上，故 $i_\infty^*j_*\cL_\psi=0$ 恰说明 $j_!\cL_\psi\to j_*\cL_\psi$ 在 $\infty$ 处的茎也为零，两者遂同构。这就是「洁净」的全部内容：边界茎消失，$!$ 延拓与 $*$ 延拓重合。零延拓那一侧的 $i_\infty^!j_!$ 恒为零，所以只需检查 $*$ 一侧。
[来源：nLab "perverse sheaf"，https://ncatlab.org/nlab/show/perverse+sheaf（Last revised 2026-10-16；2026-11-17 访问）；并见 [KL85] 与 [Gai18] 对洁净延拓的处理。]
\end{transnote}
\section{从 Satake 核到挠层}\label{sec:4.2}\index{挠层|hyperpage}
对 $V\in\Rep(\checkG)$，Satake 在局部球面仿射 Grassmann 簇上给出反常层 $\operatorname{Sat}(V)$。正则表示 $\cO(\checkG)$ 是 $\Rep(\checkG)$ 中的 ind-对象；形式上对它应用 Satake，得到一个 ind-球面 Hecke 核，而不是单个有限型反常层。对于 $G=\GL_n$，这个核是分类三元组 $(M,M',\beta:M\simeq M'\text{ 在一般点处})$ 的整体 Hecke 叠的局部模型。阿代尔双商 $\GL_n(\cO)\backslash\GL_n(\A)/\GL_n(\cO)$ 只记录选取一般位置平凡化之后的相对阿代尔位置，因此，缺少这一附加选择时，不能将它等同于完整的整体 Hecke 叠。

对 $G=\GL_n$，取出正轨迹
\[\Gr_{n,x}^+=\bigcup_{\lambda_1\ge\cdots\ge\lambda_n\ge0}\Gr_{n,x}^\lambda=(\operatorname{Mat}_{n\times n}(\cO_x)\cap\GL_n(K_x))/\GL_n(\cO_x).\]
它由在单位圆盘上定义且具有挠余核的修改组成。允许任意中心移位后，从这个正部分得到完整的仿射 Grassmann 簇。特定移位
\[(\cO_x^n\xrightarrow{\alpha}\mathcal E)\longmapsto(\cO_x(-x)^n\hookrightarrow\cO_x^n\xrightarrow{\alpha}\mathcal E)\]
使余权增加 $(1,\ldots,1)$。在整体正 Hecke 叠上，有映射
\[p:\Hecke_{\GL_n}^+\longrightarrow\Tor,\qquad(M\hookrightarrow M')\longmapsto M'/M,\]
其中 $\Hecke_{\GL_n}^+$ 分类具有挠商的秩 $n$ 丛嵌入。因此，泛正 Hecke 核应当是 $\Tor$ 上一个层的拉回。这个层就是 $\cL_E$。
\section{挠层叠}\label{sec:4.3}
定义
\[\Tor(S):=\{M\in\Coh(X\times S)\mid M\text{ 在 }S\text{ 上平坦，每个几何纤维均为挠层}\}.\]
它按长度分解：
\[\Tor=\coprod_{d\ge0}\Tor^d.\]
每个 $\Tor^d$ 都是代数叠。当 $X$ 真且 $H^0(X,\cO_X)=k$ 时，可以从 Quot 概形的开部分取得一个方便的图册：
\[\operatorname{Quot}_{\cO_X^{\oplus d}}^{\circ,d}=\{\cO_X^{\oplus d}\twoheadrightarrow M\mid M\in\Tor^d,\ H^0(\cO_X^{\oplus d})\xrightarrow{\sim}H^0(M)\}.\]
改变 $H^0(M)$ 中所选基，得到 $\GL_d$ 作用，相应商叠为
\[[\operatorname{Quot}_{\cO_X^{\oplus d}}^{\circ,d}/\GL_d]\simeq\Tor^d.\]
\srcpage{28}
\begin{example}
若 $X=\A^1$，长度 $d$ 的挠层等同于一个带有一个自同态的 $d$ 维向量空间，因此
\[\Tor^d\simeq\mathfrak{gl}_d/\GL_d.\]
\end{example}
\section{按分拆的整体分层}\label{sec:4.4}\index{分拆|hyperpage}\index{挠层!整体分层|hyperpage}
令
\[\operatorname{Part}^d=\{(d_1\ge d_2\ge\cdots\ge d_s>0)\mid d_1+\cdots+d_s=d\}.\]
叠 $\Tor^d$ 按局部闭子叠 $\Tor^{(d_1,\ldots,d_s)}$ 分层。置 $d_{s+1}=0$。由 $(d_1,\ldots,d_s)$ 标记的层是映射
\[X^{(d_1-d_2)}\times X^{(d_2-d_3)}\times\cdots\times X^{(d_s)}\longrightarrow\Tor^d\]
的像，其中
\[(D_1,\ldots,D_s)\longmapsto\cO_{D_1+\cdots+D_s}\oplus\cO_{D_2+\cdots+D_s}\oplus\cdots\oplus\cO_{D_s}.\]
分层 $\Tor^{(d)}$ 是开的；它是
\[X^{(d)}\longrightarrow\Tor^d,\qquad D\longmapsto\cO_D\]
的像。对 $X=\A^1$，这意味着
\[\Tor^{(d)}\simeq\mathfrak{gl}_d^{\mathrm{reg}}/\GL_d,\qquad\Tor^{(d)}|_{X^{(d)}-\Delta}\simeq\mathfrak{gl}_d^{\mathrm{rss}}/\GL_d,\qquad X^{(d)}\simeq\mathfrak t_d/W.\]
因此，挠层叠是在曲线情形下对 $\mathfrak{gl}_d$ 的伴随商图景的类比。
\section{\texorpdfstring{$\Tor^d$}{Torᵈ} 的 Springer 解消}\label{sec:4.5}
令 $\widetilde{\Tor}^d$ 为长度 $d$ 挠层中完整旗的叠：
\[\widetilde{\Tor}^d=\{0=M_0\subset M_1\subset\cdots\subset M_d=M\mid\length(M_i/M_{i-1})=1\}.\]
有自然映射 $\pi:\widetilde{\Tor}^d\to\Tor^d$ 及 $o_d:\widetilde{\Tor}^d\to X^d$，其中 $\pi$ 忘去旗，$o_d$ 记录商 $M_i/M_{i-1}$ 的有序支撑。记 $s_d:\Tor^d\to X^{(d)}$ 和 $r_d:X^d\to X^{(d)}$ 分别为支撑映射与对称化映射。若 $q:\Tor^{(d)}\to X^{(d)}$ 为循环层上的支撑映射，则图为
\origdiagram{p028-diagram.pdf}{0.2413}
在具有无重数支撑的正则半单轨迹 $\Tor^{d,\mathrm{rss}}$ 上，$s_d$ 将 $\Tor^{d,\mathrm{rss}}$ 展示为 $X^{(d),\mathrm{rss}}$ 上的标量自同构胚叠；$\widetilde{\Tor}^{d,\mathrm{rss}}$ 则是 $X^{d,\mathrm{rss}}$ 上的有序标量自同构胚叠。将这些标量自同构刚性化后，$\pi$ 成为通常的有序 $S_d$ 覆盖 $X^{d,\mathrm{rss}}\to X^{(d),\mathrm{rss}}$。这与 Grothendieck–Springer 图平行\cite{GPR20,AHJR14}：

\begin{transnote}{背景直觉}
「标量自同构胚叠」记录的是每条支撑所带的 $\Gm$ 自同构：无重数支撑的点 $\{x_1,\ldots,x_d\}$ 上，挠层在每个 $x_i$ 处可乘一个标量而保持同构类，这组标量构成 $\Gm^d$ 的带；胚叠的 band 就是它，所以它不是 $S_d$-覆盖而比它多出一层 $\Gm$ 的规范群。刚性化（rigidification）的做法是沿这个 band 把对象再添一层「标记」以消掉多余自同构，压掉 $\Gm^d$ 之后剩下的恰好是 $S_d$-挠子，于是 $\pi$ 变成通常的有序覆盖 $X^{d,\mathrm{rss}}\to X^{(d),\mathrm{rss}}$。这套「先把 $\Gm$-自主性记进胚叠、再刚性化掉」的手法是 [GPR20] 处理 Grothendieck–Springer 图的标准步骤。
[来源：nLab "gerbe"，https://ncatlab.org/nlab/show/gerbe（Last revised 2026-02-14；2026-11-17 访问）；并见 [GPR20] 与 [AHJR14] 对 Grothendieck–Springer 图中胚叠与刚性化的用法。]
\end{transnote}
\[\mathfrak c_d\xrightarrow{s}\mathfrak{gl}_d^{\mathrm{reg}}/\GL_d\longrightarrow\mathfrak c_d,\]
\[\widetilde{\mathfrak{gl}}_d/\GL_d\longrightarrow\mathfrak{gl}_d/\GL_d\longrightarrow\mathfrak c_d,\qquad\mathfrak t_d\longrightarrow\mathfrak c_d=\mathfrak t_d/W,\]
其中 $s$ 是 Kostant 截面。局部于 $X$，这一类比是字面上的：一点处的挠层就是幂零自同态，挠层的子层旗就是 Springer 旗。
\section{Springer–Laumon 层与 Laumon 直和项}\label{sec:4.6}
\srcpage{29}
定义次数 $d$ 的 Springer–Laumon 层为
\[\operatorname{Spr}_E^d:=\pi_!o_d^*(E^{\boxtimes d})[d]\in D_c^b(\Tor^d).\]
移位 $[d]$ 是这里等变反常约定所采用的归一化。$S_d$ 作用是 Springer 单值作用\cite{AHJR14}：在无重数支撑轨迹上，刚性化后 $\pi$ 是有序支撑的 $S_d$ 覆盖，所得作用通过中间延拓扩展到 $\operatorname{Spr}_E^d$；等价地，也可通过通常的简单反射对应得到。它不是通过任意置换一条嵌套旗而在整个旗叠上给出的字面意义上的 $S_d$ 作用。
\begin{theorem}{Laumon 延拓定理\cite{KL85}}{4.6.1}
令 $j_{\mathrm{rss}}:\Tor^{d,\mathrm{rss}}\hookrightarrow\Tor^d$ 为正则半单轨迹，$s_{\mathrm{rss}}:\Tor^{d,\mathrm{rss}}\to X^{(d),\mathrm{rss}}$ 为支撑映射，$r_{\mathrm{rss}}:X^{d,\mathrm{rss}}\to X^{(d),\mathrm{rss}}$ 为有序覆盖。层 $\operatorname{Spr}_E^d$ 是反常层，且是从正则半单轨迹作出的中间延拓：
\[\operatorname{Spr}_E^d\simeq j_{\mathrm{rss},!*}(s_{\mathrm{rss}}^*r_{\mathrm{rss},*}(E^{\boxtimes d})[d]).\]
较大的循环正则轨迹 $\Tor^{(d)}$ 也确定同一个对象；它在该轨迹上的限制，是在 $X^{(d)}$ 内跨越对角线作中间延拓所得的对象。
\end{theorem}

Laumon 层是这一 $S_d$ 作用的平凡等型直和项：
\[\cL_E^d:=\Hom_{S_d}(\mathrm{triv},\operatorname{Spr}_E^d)=(\operatorname{Spr}_E^d)^{S_d}\in\operatorname{Perv}(\Tor^d).\]
等价地，若在 $X^{(d),\mathrm{rss}}$ 上令 $E^{(d)}:=(r_{\mathrm{rss},*}E^{\boxtimes d})^{S_d}$，则
\[\cL_E^d\simeq j_{\mathrm{rss},!*}(s_{\mathrm{rss}}^*E^{(d)}[d]).\]
这个局部系统在无重数轨迹上秩为 $n^d$，不应与 $\Sym^dE$ 混淆。
\begin{lemma}{}{4.6.2}
若 $E$ 几何不可约，则 $\cL_E^d\in\operatorname{Perv}(\Tor^d)$ 几何不可约。
\end{lemma}
\begin{proof}\textbf{概要。}
基变换到 $\bar k$ 后，Schur 引理给出 $\End_{\pi_1(X_{\bar k})}(E_{\bar x})=e$，其中 $e$ 是系数域。无重数配置轨迹上的局部系统 $E^{(d)}$，对应于不可约外张量积 $E_{\bar x}^{\otimes d}$ 作为花环积群 $\pi_1(X_{\bar k})^d\rtimes S_d$ 的表示。它的自同态代数为
\[\End_{\pi_1(X_{\bar k})^d\rtimes S_d}(E_{\bar x}^{\otimes d})=(e^{\otimes d})^{S_d}=e.\]
因此，它不可约。拉回到标量自同构胚叠保持不可约性，中间延拓保持底层反常层的不可约性。
\end{proof}

最后，对 $\cL_E\in\operatorname{Perv}(\Tor)$ 定义
\[\cL_E^d:=\cL_E|_{\Tor^d}.\]
以下小节局部分析这个对象：限制到支撑于单一点的挠层后，普通 Springer 理论将局部反常对象 $\cL_{E,x}^d[-d]$ 分解为幂零层的 IC 层，其系数是 $E_x$ 的最高权表示。
\section{归约到单点处的局部挠层}\label{sec:4.7}
\srcpage{30}
为进行局部计算，以 $x$ 表示 $X$ 的一个几何点。令 $i_x:\Tor_x^d\hookrightarrow\Tor^d$ 为 $s_d$ 在 $d\cdot x$ 上的纤维。因此，$\Tor_x^d$ 分类支撑于 $x$ 的长度 $d$ 挠层。置
\[\cL_{E,x}^d:=i_x^*\cL_E^d.\]
由于这一纤维在对称幂方向上余维为 $d$，局部反常对象是 $\cL_{E,x}^d[-d]$。
\begin{lemma}{不交支撑的因子化}{4.7.1}
若 $x\ne y$，直和给出局部挠层叠的同构
\[\Tor_x^a\times\Tor_y^b\xrightarrow{\sim}\Tor_{ax+by}^{a+b},\qquad\cL_{E,x}^a\boxtimes\cL_{E,y}^b\xrightarrow{\sim}\cL_{E,ax+by}^{a+b}.\]
更一般地，对具有不同支撑点的除子 $D=\sum_xd_xx$，
\[\prod_x\Tor_x^{d_x}\xrightarrow{\sim}\Tor_D^d,\qquad\boxtimes_x\cL_{E,x}^{d_x}\simeq\cL_{E,D}^d.\]
\end{lemma}

因此，$\cL_E^d$ 的结构归约为计算 $\cL_{E,x}^d[-d]\in\operatorname{Perv}(\Tor_x^d)$。
\section{按分拆的局部分层}\label{sec:4.8}\index{挠层!局部分层|hyperpage}
设 $\lambda=(\lambda_1\ge\lambda_2\ge\cdots\ge\lambda_s>0)$ 为 $d$ 的分拆，置 $\lambda_{s+1}=0$。记 $|\lambda|=d$，$\ell(\lambda)=s$。对有效除子 $D$，记 $\cO_D:=\cO_X/\cO_X(-D)$。整体分层 $\Tor^\lambda\subset\Tor^d$ 由图册
\[X^{(\lambda_1-\lambda_2)}\times X^{(\lambda_2-\lambda_3)}\times\cdots\times X^{(\lambda_s)}\longrightarrow\Tor^\lambda\]
描述，它将 $(D_1,\ldots,D_s)$ 映为
\[\cO_{D_1+\cdots+D_s}\oplus\cO_{D_2+\cdots+D_s}\oplus\cdots\oplus\cO_{D_s}.\]
分层 $\Tor^\lambda$ 局部闭地映入 $\Tor^d$。

$x$ 处的局部层是剩余部分 $\Tor_x^\lambda\subset\Tor_x^d$，对应于挠层
\[T_{\lambda,x}:=\bigoplus_{i=1}^s\cO_X/\cO_X(-\lambda_ix).\]
等价地，过渡到几何点 $x$ 处的完备局部环并选取素元 $t$ 后，它是 $k[[t]]$ 模
\[\bigoplus_{i=1}^sk[[t]]/(t^{\lambda_i}).\]
叠 $\Tor_x^d$ 的维数为 $-d$；在下述幂零锥模型中，它的各分层是幂零轨道的商叠。

\begin{transnote}{含糊点拨}
维数可以算出来：幂零锥 $\cN_d\subset\mathfrak{gl}_d$ 由 $\det=0$（一个方程）在 $d^2$ 维的 $\mathfrak{gl}_d$ 中切出，故 $\dim\cN_d=d^2-d$；商去 $\GL_d$ 的 $d^2$ 维轨道，得 $\dim\Tor_x^d=(d^2-d)-d^2=-d$。负维数不是笔误：商叠的维数按「维数减群维数」计，正是「$\GL_d$ 的作用比底空间还大」的表现。各分层 $\Tor_x^\lambda$ 则是幂零轨道 $\cO_\lambda$ 的商叠，维数同样要减去 $d^2$——这与 \S\ref{sec:4.8} 末尾的幂零锥模型 $\Tor_x^d\simeq\cN_d/\GL_d$ 一致。
[来源：nLab "algebraic stack"，https://ncatlab.org/nlab/show/algebraic+stack（Last revised 2023-04-16；2026-11-17 访问）；并见 [GPR20] 对 $\Tor^d$ 的幂零锥描述（原书 p.30）。]
\end{transnote}
\begin{definition}{}{4.8.1}
对分拆 $\lambda\vdash d$，定义
\[B_{\lambda,x}:=\operatorname{IC}_{\overline{\Tor_x^\lambda}}.\]
\end{definition}
\section{Laumon 层的局部分解}\label{sec:4.9}
令 $V=E_x$ 为 $E$ 在 $x$ 处的纤维。对满足 $s\le n$ 的分拆 $\lambda=(\lambda_1\ge\cdots\ge\lambda_s>0)$，添零将其扩展为 $\GL(V)\simeq\GL_n$ 的优势权：
\[\lambda=(\lambda_1,\ldots,\lambda_s,0,\ldots,0).\]
令 $V^\lambda$ 为 $\GL(V)$ 中最高权为 $\lambda$ 的不可约最高权表示。

等价地，若 $\chi^\lambda\in\operatorname{Irr}(S_d)$ 是与 $\lambda$ 对应的 Specht 表示，则 Schur–Weyl 对偶给出
\[V^\lambda\simeq(V^{\otimes d}\otimes\chi^\lambda)^{S_d}.\]
若 $s>n$，这个向量空间为零。
\srcpage{31}
\begin{proposition}{局部公式}{4.9.1}
对每个几何点 $x\in X$ 与整数 $d\ge0$，置 $n(\lambda):=\sum_{i\ge1}(i-1)\lambda_i$。按通常的 Casselman–Shalika 归一化，将其看作 $n(\lambda)=\langle\lambda,(0,1,\ldots,n-1)\rangle$。则
\[\cL_{E,x}^d[-d]\simeq\bigoplus_{\substack{\lambda\vdash d\\\ell(\lambda)\le n}}B_{\lambda,x}\otimes E_x^\lambda(-n(\lambda)).\]
\end{proposition}
\section{与仿射 Grassmann 簇的关系}\label{sec:4.10}
令 $\Gr_{n,x}$ 为 $\GL_n$ 在 $x$ 处的仿射 Grassmann 簇。令 $\Gr_{n,x}^{d,\mathrm{mod}}$ 为次数 $d$ 的正修改轨迹，即其点是 $\cO_x^n$ 的长度 $d$ 商的轨迹；这一记号不应与正则幂零轨迹混淆。令
\[f:\Gr_{n,x}^{d,\mathrm{mod}}\longrightarrow\Tor_x^d\]
为将修改映为其商挠层的映射。它的像是最多可由 $n$ 个元素生成的挠层所构成的开子叠；在这一像上，$f$ 光滑，相对维数为 $nd$。对 $|\lambda|=d$ 且 $\ell(\lambda)\le n$，令 $\Gr_{n,x}^\lambda$ 为型 $\lambda$ 的 Schubert 分层。图为
\origdiagram{p031-square.pdf}{0.2596}
置
\[A_\lambda:=\operatorname{IC}_{\overline{\Gr_{n,x}^\lambda}}.\]
采用 IC 包括反常上同调移位但不包括 Tate 扭转的约定，光滑拉回给出
\[f^*B_{\lambda,x}\simeq A_\lambda[-dn].\]
等价地，$f^*B_{\lambda,x}[dn]\simeq A_\lambda$。局部 Laumon 公式中的 Tate 因子 $-n(\lambda)$ 属于 Whittaker／Casselman–Shalika 归一化，而不属于 IC 层的光滑拉回比较。
\newpage
\section{用 Springer 理论证明局部公式}\label{sec:4.11}
\begin{proof}
现在使用经典 Springer 对应\cite{GPR20,AHJR14}说明局部分解的证明。在 $x$ 的 étale 局部工作，或将 $X$ 替换为 $x$ 处的形式圆盘。选取素元后，支撑于 $x$ 的长度 $d$ 挠层就是一个带有幂零自同态的 $d$ 维向量空间。因此
\[\Tor_x^d\simeq\cN_d/\GL_d,\]
其中 $\cN_d\subset\mathfrak{gl}_d$ 是幂零锥。非局部陈述 $\Tor^d\simeq\mathfrak{gl}_d/\GL_d$ 对 $X=\A^1$ 是字面成立的；对一般曲线，它是在次数 $d$ 除子附近相应的 étale 局部模型。在这一局部等同下，映射 $\widetilde{\Tor}^d\to\Tor^d$ 成为 Grothendieck–Springer 解消，限制到 $\Tor_x^d$ 上则成为幂零锥的 Springer 解消。

基变换到几何点 $\bar x$，并暂时省去 Tate 扭转。局部反常 Springer 层分解为
\[\operatorname{Spr}_{\bar x}^d[-d]\simeq\bigoplus_{\chi\in\operatorname{Irr}(S_d)}\chi\otimes\operatorname{IC}_{\overline{\cO_\chi}}.\]
\srcpage{32}
这里，$\operatorname{Irr}(S_d)$ 等同于 $d$ 的分拆，也等同于 $\mathfrak{gl}_d$ 中的幂零轨道。采用以下约定：
\[\lambda\longleftrightarrow\cO_\lambda,\qquad\text{平凡表示}\longleftrightarrow\text{正则幂零轨道 }(d).\]
这确定了不取转置的约定。

对 Laumon 层，纤维 $E_x^{\otimes d}$ 随 Springer 层一起带入，$S_d$ 通过置换 $d$ 个因子作用于这个张量因子。因此
\[\operatorname{Spr}_{E,x}^d[-d]\simeq\bigoplus_{\chi\in\operatorname{Irr}(S_d)}\chi\otimes E_x^{\otimes d}\otimes\operatorname{IC}_{\overline{\cO_\chi}}.\]
取 $S_d$ 不变直和项得到
\[(\operatorname{Spr}_{E,x}^d[-d])^{S_d}\simeq\bigoplus_{\chi\in\operatorname{Irr}(S_d)}(\chi\otimes E_x^{\otimes d})^{S_d}\otimes\operatorname{IC}_{\overline{\cO_\chi}}.\]
现在写 $\chi=\chi^\lambda$，其中 $d=\lambda_1+\cdots+\lambda_s$、$\lambda_1\ge\cdots\ge\lambda_s>0$。添零将 $\lambda$ 扩展到长度 $n$。由 Schur–Weyl 对偶，
\[(\chi^\lambda\otimes E_x^{\otimes d})^{S_d}\simeq E_x^\lambda,\]
若 $s>\dim(E_x)=n$ 则它消失。最后，在 $\Tor_x^d\simeq\cN_d/\GL_d$ 下，有
\[\operatorname{IC}_{\overline{\cO_\lambda}}=B_{\lambda,x}.\]
恢复 Tate 扭转，便得到命题\ref{pro:4.9.1}关于 $\cL_{E,x}^d[-d]$ 的公式。
\end{proof}
